\documentclass[11pt,a4paper]{article}
\usepackage[margin=2.2cm]{geometry}
\usepackage{fontspec}
\IfFontExistsTF{Avenir Next}{
  \setmainfont{Avenir Next}
}{
  \IfFontExistsTF{Helvetica Neue}{
    \setmainfont{Helvetica Neue}
  }{
    \IfFontExistsTF{Source Sans 3}{
      \setmainfont{Source Sans 3}
    }{
      \setmainfont{Times New Roman}
    }
  }
}
\usepackage[hidelinks]{hyperref}
\setlength{\parindent}{0pt}
\setlength{\parskip}{0.8em}

\begin{document}

{\LARGE \textbf{Cristi Prodan}}\\
Telefon: +40 746 708 927\\
Email: \href{mailto:prodan.cristian@gmail.com}{prodan.cristian@gmail.com}\\
GitHub: \href{https://github.com/christian}{https://github.com/christian}\\
LinkedIn: \href{https://www.linkedin.com/in/cristianprodan/}{https://www.linkedin.com/in/cristianprodan/}

\vspace{1.2em}

{\Large \textbf{Scrisoare de intenție}}

Bună ziua,

Vă trimit candidatura mea pentru rolul din echipa dumneavoastră, deoarece experiența mea se potrivește foarte bine cu responsabilitățile de infrastructură cloud, operare de servicii critice și colaborare cu echipe tehnice multiple.

Am peste 10 ani de experiență în dezvoltarea, operarea și mentenanța serviciilor cloud și backend, în special în medii AWS și GCP, cu responsabilități directe pentru microservicii, infrastructură, CI/CD, securitate, observabilitate și răspuns la incidente. În rolul meu de Senior Software Engineer la PTC Schweiz, am construit, deținut și operat servicii Vuforia Cloud în producție, am lucrat cu Kubernetes, CloudFormation, Terraform, GitHub Actions și TeamCity, și am contribuit la scalare, toleranță la erori, load balancing, securitate și monitorizare cu Elasticsearch, Logstash și Kibana.

Deși experiența mea principală nu este centrată pe Microsoft Azure, profilul meu tehnic este foarte apropiat de cerințele rolului: am lucrat cu servicii serverless, infrastructură cloud, pipeline-uri CI/CD, autentificare SSO, API-uri de producție și sisteme distribuite. Am experiență practică în proiectarea arhitecturilor, documentarea API-urilor cu OpenAPI, integrarea serviciilor externe și investigarea incidentelor în medii cloud. Sunt obișnuit să intru rapid în ecosisteme noi și să le duc în producție într-un mod disciplinat și fiabil.

Cred că pot aduce valoare echipei prin combinația dintre experiența hands-on în cloud, gândirea de sistem, atenția la operare și capacitatea de a colabora cu echipe centrale, echipe de produs, testare și suport. Îmi place să lucrez cu standarde clare, să documentez deciziile tehnice și să construiesc soluții care pot fi operate și întreținute pe termen lung.

Aș aprecia ocazia să discutăm mai concret despre nevoile echipei și despre cum experiența mea în infrastructură cloud, servicii backend și operare de producție poate contribui la proiectele dumneavoastră.

Cu stimă,\\
Cristi Prodan

\end{document}
